\section{总结与延伸}

\begin{enumerate}
    \item 探索是 RL 的核心挑战，$\epsilon$-greedy 到 UCB/Thompson Sampling 形成了一个日益精细的工具链
    \item 深度 RL 中需要用密度模型或网络预测误差来近似新颖性
    \item 基于好奇心的方法需警惕噪声电视问题
    \item 元探索将"如何探索"本身作为可学习的能力
    \item 在机器人和 LLM 等真实应用中，有效探索直接决定了学习效率
\end{enumerate}

\subsection{应用外延：把 meta-exploration 用到计算机教育}

课程最后还展示了一个很少见但很有启发性的方向：把 meta-RL 的探索策略用于自动批改和学生程序反馈。系统先学习 ``该运行哪些测试、该看哪些行为信号最有信息量''，再更快地找出 bug、生成评分建议。

\newpage
\begin{figure}[H]
\centering
\includegraphics[width=0.9\textwidth]{slides-images/slide-30.jpg}
\caption{Meta-exploration 在程序反馈中的应用：探索哪些交互最能暴露学生代码问题}
\end{figure}

\subsection{教育场景里的 ``探索'' 和传统 RL 有何不同}

这里的探索不再是操控机器人或玩游戏，而是决定：\textbf{为了判断学生程序哪里有问题，系统应该优先运行哪些交互、观察哪些信号、展示哪些视频片段。} 这类探索的价值在于节省助教时间，同时让反馈更聚焦、更可操作。

\begin{figure}[H]
\centering
\includegraphics[width=0.9\textwidth]{slides-images/slide-31.jpg}
\caption{Bounce 作业中的 learned exploration：不同交互会暴露不同类别的错误}
\end{figure}

\begin{table}[H]
\centering
\begin{tabular}{p{0.2\textwidth} p{0.28\textwidth} p{0.28\textwidth}}
\toprule
教育任务 & 需要探索的信息 & 系统收益 \\
\midrule
交互式编程作业反馈 & 哪些测试情境最能触发 bug & 更快定位错误来源 \\
助教评分辅助 & 哪些证据最足以支持 rubric 打分 & 减少人工查看和重复判断 \\
错误模式分析 & 哪些学生行为属于同一类失误 & 形成更可复用的反馈模板 \\
\bottomrule
\end{tabular}
\caption{Meta-exploration 在 CS 教育中的信息收集目标}
\end{table}

\begin{figure}[H]
\centering
\includegraphics[width=0.9\textwidth]{slides-images/slide-32.jpg}
\caption{AI-assisted grading 的课程实践：更快且更准确的助教评分流程}
\end{figure}

\subsection{本讲总表}

\begin{table}[H]
\centering
\begin{tabular}{p{0.18\textwidth} p{0.22\textwidth} p{0.22\textwidth} p{0.22\textwidth}}
\toprule
路线 & 核心思想 & 优势 & 主要问题 \\
\midrule
$\epsilon$-greedy / UCB / Thompson & 在 bandit 中用不确定性指导探索 & 可分析、可证明 regret & 难直接扩展到大规模 MDP \\
Deep RL intrinsic exploration & 用新颖性或预测误差鼓励探索 & 适配连续高维状态空间 & 易被噪声和伪新颖性误导 \\
End-to-end meta-exploration & 直接从任务 reward 学探索与执行 & 原理统一，理论上最优 & coupling 问题导致训练困难 \\
DREAM-style decoupling & 先学任务关键信息，再学恢复它的探索 & 更易优化，实践效果好 & 依赖任务表示与信息瓶颈设计 \\
\bottomrule
\end{tabular}
\caption{Lecture 14 的探索方法地图}
\end{table}

\begin{knowledgebox}{拓展阅读}
\begin{itemize}
    \item Auer et al., ``Finite-time Analysis of the Multiarmed Bandit Problem,'' Machine Learning 2002
    \item Bellemare et al., ``Unifying Count-Based Exploration and Intrinsic Motivation,'' NeurIPS 2016
    \item Burda et al., ``Exploration by Random Network Distillation,'' ICLR 2019
    \item Pathak et al., ``Curiosity-driven Exploration by Self-Supervised Prediction,'' ICML 2017
    \item Liu et al., ``Explore then Execute: Adapting without Rewards via Factorized Meta-Reinforcement Learning,'' 2022
\end{itemize}
\end{knowledgebox}

\end{document}
